% FIG ch10-02 paper-to-card (Chinese)
\usetikzlibrary{arrows.meta,calc,fit,positioning}
\begin{tikzpicture}[
  font=\sffamily\footnotesize,
  node distance=4mm and 6mm,
  box/.style={
    draw=BookRule, line width=0.4pt, rounded corners=2pt,
    align=center, inner sep=5pt, text=BookInk, text width=22mm,
  },
  autobox/.style={box, dashed, fill=BookPale!40},
  authorbox/.style={box, fill=white},
  limitbox/.style={authorbox, draw=BookTeal, line width=0.9pt},
  flow/.style={-{Stealth[length=2mm]}, draw=BookInk!55, thin},
  note/.style={font=\scriptsize, text=BookInk!75, align=left},
]
  \node[authorbox, minimum width=24mm, minimum height=28mm] (pdf) {%
    \textbf{论文 PDF}\\[2pt]
    \tikz{\draw[fill=BookPale!50, draw=BookInk!30] (0,0) rectangle (1.4,1.8);
      \draw[fill=BookTeal!25, draw=BookTeal!60] (0,0.9) rectangle (1.4,1.8);
      \node[font=\tiny, anchor=west] at (1.55,1.35) {前 5 页};
      \node[font=\tiny, anchor=west] at (1.55,0.45) {其余页};}};
  \node[autobox, right=of pdf] (auto) {PDF\\转提示词};
  \node[autobox, right=of auto, text width=28mm] (prompt) {60--110 词的\\视觉提示词};
  \node[authorbox, below=10mm of pdf, text width=26mm] (full) {通读\\整篇论文};
  \node[authorbox, right=of full, text width=30mm] (card) {%
    贡献卡四行\\[1pt]
    \begin{tabular}{@{}c@{}}
      问题\\方法\\发现\\\textbf{限定}
    \end{tabular}};
  \node[limitbox, fit=(card), inner sep=1pt] {};
  \node[authorbox, right=12mm of prompt, text width=32mm] (merge) {逐行比对：\\补回限定条件、\\删去超出论文的说法};
  \node[authorbox, right=of merge, text width=28mm] (final) {作者确认的\\生成提示词};
  \draw[flow] (pdf) -- (auto);
  \draw[flow] (auto) -- (prompt);
  \draw[flow] (full) -- (card);
  \draw[flow] (prompt.east) -| (merge.west);
  \draw[flow] (card.east) -| (merge.south);
  \draw[flow] (merge) -- (final);
  \node[note, anchor=north west] at (pdf.south west) {约 6{,}000 字符上限};
  \node[note, anchor=north] at ($(auto)!0.5!(prompt)+(0,-8mm)$) {自动路径（虚线框）};
  \node[note, anchor=north] at ($(full)!0.5!(card)+(0,-8mm)$) {作者路径（实线框）};
  \node[note, anchor=north east] at (final.south east) {向导与工具不检查科学内容};
\end{tikzpicture}
